The exploratory record \(\mathcal J_{\mathrm{PICC}}\) is the reproducible but unexecuted procedure implemented by \(\texttt{discovery/jiangxi\_robust\_reanalysis.py}\).  Its input directory must contain the Dataverse-converted tables \(\texttt{survey.tab}\) and \(\texttt{network.tab}\), which are the filenames the script actually consumes.  The procedure cleans identifiers, sentinel \(99\), self-loops, and duplicate directed nominations while retaining cross-village nominations, constructs the three target populations for \(\tau_{10},\tau_{21},\tau_{32}\), and uses the disclosed independent Bernoulli \((1/2)\) reconstruction necessitated by the unreleased original strata.  It splits connected components of the exposure-overlap graph without cutting components, estimates the two marginal laws and Fr\'echet endpoints on each training fold, and optimizes robust and zero-covariance bases under the same initialization list and iteration budget.  Each held-out basis is deployed through def:observable-coefficient-map with Zhao's ordinary inverse and the additional \(\Pi^{-1}\) factor; no Moore--Penrose safeguard is used or attributed to Zhao.  The output fields report point estimates, design-randomization standard errors, both endpoint losses, worst loss, minimum regularity, and a heuristic objective diagnostic.  The latter is not a certified upper-minus-lower optimization gap.  The authenticated replication data were unavailable in the solve environment, so the protocol was not executed and no numerical table is asserted.  It remains exploratory because its internal cross-fitting is not the independent unpaired pilot of ass:unpaired-pilot.